% =========================================================================
% ACTA DE REUNIÓN — CS3081-004-2026 · BORRADOR
% Plantilla de estilo basada en el modelo oficial de Google Docs del curso
% Compilar con:  tectonic Acta_Reunion4_BORRADOR.tex
% =========================================================================
\documentclass[11pt,a4paper]{article}

\usepackage[a4paper, top=2.6cm, bottom=2.4cm, left=2.5cm, right=2.5cm]{geometry}
\usepackage{fontspec}
\setmainfont{Carlito}                    % métrica idéntica a Calibri
\usepackage[table]{xcolor}
\usepackage{array}
\usepackage{longtable}
\usepackage{fancyhdr}
\usepackage{lastpage}
\usepackage{needspace}

% ------------------------- Paleta de la plantilla -------------------------
\definecolor{rojo}{HTML}{E4002B}   % marca CS3081
\definecolor{grisB}{HTML}{B7B7B7}  % bordes de tabla
\definecolor{grisH}{HTML}{202124}  % encabezados de tabla (negro)
\definecolor{grisL}{HTML}{F3F4F6}  % celdas de etiqueta
\definecolor{rosa}{HTML}{FCE8EC}   % caja PROPÓSITO / REGISTRO
\definecolor{grisV}{HTML}{555555}  % valores (itálica)
\definecolor{grisN}{HTML}{666666}  % notas y pie
\definecolor{tinto}{HTML}{333333}

\arrayrulecolor{grisB}
\renewcommand{\arraystretch}{1.55}
\setlength{\tabcolsep}{4pt}

% ------------------------- Encabezado y pie corridos ----------------------
\pagestyle{fancy}
\fancyhf{}
\fancyhead[L]{\footnotesize\textbf{\color{rojo}CS3081\hspace{0.7em}|\hspace{0.7em}INGENIERÍA DE SOFTWARE}\hspace{1.4em}{\color{grisN}·\hspace{0.7em}Acta de reunión con usuarios}}
\fancyfoot[R]{\footnotesize\color{grisN}Uso académico · Proyecto real\hspace{2em}Página \thepage\ de \pageref{LastPage}}
\renewcommand{\headrulewidth}{0pt}
\renewcommand{\footrulewidth}{0pt}
\setlength{\headheight}{14pt}

% ------------------------- Comandos de estilo -----------------------------
% Sección: 14pt, con aire vertical (no pegada al contenido).
% Needspace: si no queda espacio mínimo (~6cm) para título + nota + primeras
% filas, la sección salta COMPLETA a la página siguiente (nunca queda huérfana).
\newcommand{\asec}[1]{\Needspace*{13\baselineskip}\par\vspace{18pt}{\noindent\fontsize{14}{17}\selectfont\textbf{#1}}\par\vspace{9pt}}
% Caja de nota (PROPÓSITO / REGISTRO)
\newcommand{\cajanota}[2]{\par\noindent\fcolorbox{grisB}{rosa}{\parbox{\dimexpr\linewidth-2\fboxsep-2\fboxrule}{\footnotesize \textbf{\textcolor{rojo}{#1: }}{\color{tinto}#2}}}\par\vspace{7pt}}
% Celda encabezado de tabla / etiqueta / valor
\newcommand{\hc}[1]{\cellcolor{grisH}\textcolor{white}{\textbf{#1}}}
\newcommand{\lb}[1]{\cellcolor{grisL}\textbf{#1}}
\newcommand{\vl}[1]{\textcolor{grisV}{\textit{#1}}}
% Columnas
\newcolumntype{L}[1]{>{\raggedright\arraybackslash}p{#1}}
\newcolumntype{C}[1]{>{\centering\arraybackslash}p{#1}}

\setlength{\parindent}{0pt}
\setlength{\parskip}{4pt}

\begin{document}

% ============================ TÍTULO ======================================
{\centering {\fontsize{22}{26}\selectfont\textbf{ACTA DE REUNION}}\par}
\vspace{10pt}

\cajanota{PROPÓSITO}{documentar cada interacción con usuarios o representantes del cliente. Complete el acta durante la reunión, compártala dentro de las 24 horas y registre la evidencia de validación.}

% ============================ 1. DATOS ====================================
\asec{1. DATOS GENERALES}
{\footnotesize
\begin{tabular}{|L{2.35cm}|L{5.45cm}|L{2.1cm}|L{4.3cm}|}
\hline
\lb{Código de acta} & \vl{CS3081-004-2026} & \lb{Fecha} & \vl{08/09/2026} \\ \hline
\lb{Hora de inicio} & \vl{10:55} & \lb{Hora de término} & \vl{11:15} \\ \hline
\lb{Proyecto} & \multicolumn{3}{L{12.0cm}|}{\vl{Proyecto Igualab - Plataforma de analítica de sostenibilidad (RAG)}} \\ \hline
\lb{Módulo/Fase} & \vl{Análisis y Diseño --- confirmación de servidor y cierre del análisis (arquitectura)} & \lb{Modalidad} & \vl{[ ] Presencial \ [X] Virtual} \\ \hline
\end{tabular}}

% ============================ 2. OBJETIVO =================================
\asec{2. OBJETIVO Y RESULTADO ESPERADO}
{\footnotesize
\begin{tabular}{|L{2.9cm}|L{11.3cm}|}
\hline
\lb{Objetivo de la reunión} & \vl{Confirmar la infraestructura/servidor para la fase 1 (pendiente del acta 3), revisar el análisis GAP y el backlog inicial, presentar el Análisis y Diseño v1.1 (proceso to-be) identificando los artefactos que faltan para cerrar el análisis con arquitectura (diagramas de casos de uso y diagrama de base de datos con diccionario de datos), someter a decisión del cliente el modelo de usuarios del portal y formalizar el cambio de roles del equipo (nuevo PM).} \\ \hline
\lb{Resultado esperado} & \vl{Servidor/despliegue confirmado (fase 1 en desarrollo, sin producción), GAP y backlog revisados, pendientes de análisis comprometidos con fecha (11-09), modelo de usuarios decidido por el cliente y nuevo PM formalizado.} \\ \hline
\end{tabular}}

% ============================ 3. PARTICIPANTES ============================
\asec{3. PARTICIPANTES}
{\footnotesize
\begin{tabular}{|C{0.8cm}|L{4.1cm}|L{3.1cm}|L{4.4cm}|C{1.8cm}|}
\hline
\hc{N°} & \hc{Nombre y apellido} & \hc{Organización / equipo} & \hc{Rol en la reunión} & \hc{Asistencia} \\ \hline
1 & \vl{Oscar Rafael Baldeón Montoro} & \vl{Cliente} & \vl{Cliente / Product Owner} & \vl{[ ] Sí \ [X] No} \\ \hline
2 & \vl{\textbf{Fabricio Godofredo Ladera La Torre}} & \vl{Equipo UTEC} & \vl{\textbf{Jefe de Proyecto (PM)}} & \vl{[X] Sí \ [ ] No} \\ \hline
3 & \vl{Juan Renato Flores Pascual} & \vl{Equipo UTEC} & \vl{Desarrollador Frontend} & \vl{[ ] Sí \ [X] No} \\ \hline
4 & \vl{Carlos Ordinola} & \vl{Equipo UTEC} & \vl{Analista Funcional} & \vl{[ ] Sí \ [X] No} \\ \hline
5 & \vl{Luis Millones} & \vl{Equipo UTEC} & \vl{Analista funcional} & \vl{[ ] Sí \ [X] No} \\ \hline
6 & \vl{Juan Ferreyra} & \vl{Equipo UTEC} & \vl{Backend developer} & \vl{[ ] Sí \ [X] No} \\ \hline
7 & \vl{Mauricio Gonzales Krener} & \vl{Equipo UTEC} & \vl{Tester} & \vl{[ ] Sí \ [X] No} \\ \hline
8 & \vl{Alonso Benites Camacho} & \vl{Equipo UTEC} & \vl{Tester} & \vl{[ ] Sí \ [X] No} \\ \hline
\end{tabular}}

% ============================ 4. AGENDA ===================================
\asec{4. AGENDA}
{\footnotesize
\begin{tabular}{|C{0.8cm}|L{8.6cm}|L{2.6cm}|C{2.2cm}|}
\hline
\hc{N°} & \hc{Tema} & \hc{Responsable} & \hc{Tratado} \\ \hline
1 & \vl{Formalización del cambio de roles del equipo (nuevo PM)} & \vl{Fabricio (PM)} & \vl{[X] Sí \ [ ] No} \\ \hline
2 & \vl{Confirmación de servidor y despliegue de la fase 1 (sin producción)} & \vl{Oscar / Equipo} & \vl{[X] Sí \ [ ] No} \\ \hline
3 & \vl{Revisión del análisis GAP y backlog inicial} & \vl{Equipo / Oscar} & \vl{[X] Sí \ [ ] No} \\ \hline
4 & \vl{Análisis y Diseño v1.1 (to-be) y pendientes para cerrar el análisis con arquitectura} & \vl{Carlos / Luis} & \vl{[X] Sí \ [ ] No} \\ \hline
5 & \vl{Decisión del modelo de usuarios del portal (RBAC)} & \vl{Oscar (Cliente)} & \vl{[X] Sí \ [ ] No} \\ \hline
\end{tabular}}

% ============================ 5. DESARROLLO ===============================
\asec{5. DESARROLLO DE LA REUNIÓN}
{\footnotesize
\begin{longtable}{|C{0.8cm}|L{3.0cm}|L{7.6cm}|L{2.8cm}|}
\hline
\rowcolor{grisH}\hc{N°} & \hc{Tema / necesidad} & \hc{Síntesis de lo conversado} & \hc{Conclusión / estado} \\ \hline
\endfirsthead
\hline
\rowcolor{grisH}\hc{N°} & \hc{Tema / necesidad} & \hc{Síntesis de lo conversado} & \hc{Conclusión / estado} \\ \hline
\endhead
1 & \vl{Cambio de roles del equipo} & \vl{Juan Renato Flores Pascual pasa a desempeñarse como Desarrollador Frontend y Fabricio Godofredo Ladera La Torre asume como Jefe de Proyecto (PM). El cliente toma conocimiento del cambio.} & \vl{Decidido} \\ \hline
2 & \vl{Confirmación de servidor y despliegue} & \vl{Respuesta a lo pendiente del acta 3 (REQ-08): toda la \textbf{fase 1 operará en el entorno de desarrollo} provisto por la universidad (servidor + LLM por API --- acta 2). \textbf{No habrá despliegue a producción en fase 1}.} & \vl{Decidido (cierra REQ-08)} \\ \hline
3 & \vl{Análisis GAP y backlog inicial} & \vl{Se revisó el análisis GAP y el backlog inicial del producto. Se aprueban con los ajustes tratados en esta acta y quedan priorizados para el Sprint 1 (login y RBAC).} & \vl{Aprobado} \\ \hline
4 & \vl{Análisis y Diseño v1.1 y cierre del análisis} & \vl{Se presentó la versión actualizada del documento con el \textbf{proceso to-be}. El análisis \textbf{aún no se cierra}: faltan los \textbf{diagramas de casos de uso} y el \textbf{diagrama de base de datos con diccionario de datos}. El equipo dispone de \textbf{1--2 semanas adicionales} para el documento; su avance depende de la \textbf{confirmación del cliente} sobre el diagrama de procesos (ajuste del modelo de usuarios), punto central de esta reunión.} & \vl{En proceso (sujeto a confirmación)} \\ \hline
5 & \vl{Modelo de usuarios del portal} & \vl{El equipo planteó la duda sobre el RBAC: \textbf{(a)} 2 roles / 3 personas --- Oscar = Superadmin + 2 Administradores, sin rol ``Usuario'' (lo comentado por el gerente), o \textbf{(b)} mantener los 3 roles internos --- Superadmin = Oscar, Administrador = 2 de Oscar, Usuario (solo lectura) = personal de Igualab, lo que ajustaría el ``máx. 3 usuarios'' del acta 1 (REQ-03). \textbf{Decisión del cliente:} se adopta la \textbf{opción (a)} --- RBAC de \textbf{2 roles / 3 personas} (Oscar = Superadmin + 2 Administradores; sin rol ``Usuario''), habilitando el avance del diagrama de procesos.} & \vl{Definido --- opción (a)} \\ \hline
6 & \vl{Alcance futuro (2ª fase)} & \vl{Se registra como alcance de \textbf{2ª fase} (sin desarrollo en fase 1): alta de \textbf{usuarios ``Cliente'' con tenant\_id} (modelo multi-tenant) y evaluación de una \textbf{pasarela de pagos} para la suscripción definida en las actas 1 y 2. \textbf{Aclaración importante:} todo lo contemplado en el plan de proyecto corresponde \textbf{únicamente a la fase 1}; la fase 2 no está comprometida y dependerá del \textbf{tiempo disponible del equipo} (cronograma académico) y del \textbf{presupuesto del cliente para salir a producción} (servidor e infraestructura IA/RAG). De concretarse esa inversión, la plataforma pasaría de herramienta interna de prospección a \textbf{producto generador de ingresos} --- suscripción de empresas con acceso a sus reportes y soporte de la consultoría retribuida ---, con lo cual la propia fase 2 amortizaría el costo de producción.} & \vl{Registrado (2ª fase)} \\ \hline
\end{longtable}}

% ============================ 6. REQUISITOS ===============================
\asec{6. NECESIDADES, REQUISITOS Y CAMBIOS IDENTIFICADOS}
{\footnotesize\itshape\color{grisN} Registre solo lo comprendido en la reunión. La incorporación al alcance se confirma mediante la priorización y validación correspondiente.\par}
\vspace{5pt}
{\footnotesize
\begin{longtable}{|C{1.3cm}|L{5.3cm}|C{1.3cm}|C{1.3cm}|L{3.4cm}|C{1.6cm}|}
\hline
\rowcolor{grisH}\hc{ID} & \hc{Necesidad / requisito / cambio} & \hc{Tipo} & \hc{Prioridad} & \hc{Criterio o evidencia} & \hc{Estado} \\ \hline
\endfirsthead
\hline
\rowcolor{grisH}\hc{ID} & \hc{Necesidad / requisito / cambio} & \hc{Tipo} & \hc{Prioridad} & \hc{Criterio o evidencia} & \hc{Estado} \\ \hline
\endhead
REQ-10 & \vl{Definir el \textbf{modelo de usuarios}: (a) 2 roles / 3 personas (Oscar = Superadmin + 2 Administradores, sin rol ``Usuario'') o (b) 3 roles internos (Superadmin = Oscar · Administrador = 2 de Oscar · Usuario = personal de Igualab, solo lectura), lo que ajusta el \textbf{``máx. 3 usuarios''} del acta 1 (REQ-03).} & \vl{Cambio} & \vl{Alta} & \vl{Decisión final del cliente en la reunión} & \vl{Definido} \\ \hline
REQ-11 & \vl{La \textbf{fase 1 opera íntegramente en entorno de desarrollo} (infraestructura de la universidad: servidor + LLM por API); \textbf{sin despliegue a producción}.} & \vl{RN} & \vl{Alta} & \vl{Sin presupuesto del cliente para producción (cierra REQ-08)} & \vl{Definido} \\ \hline
REQ-12 & \vl{Alcance 2ª fase: usuarios \textbf{``Cliente'' con tenant\_id} (multi-tenant) + \textbf{pasarela de pagos} (a evaluar) para el modelo de suscripción. \textbf{Aclaración:} todo lo contemplado en el plan de proyecto corresponde \textbf{únicamente a la fase 1}.} & \vl{Cambio} & \vl{Media} & \vl{Registrado para 2ª fase; fuera del desarrollo de fase 1} & \vl{Definido} \\ \hline
REQ-13 & \vl{\textbf{Cerrar el análisis con arquitectura}: elaborar los \textbf{diagramas de casos de uso} y el \textbf{diagrama de base de datos con diccionario de datos}.} & \vl{Documento} & \vl{Alta} & \vl{Estimado en 1--2 semanas tras la confirmación del cliente (diagrama de procesos)} & \vl{En proceso} \\ \hline
REQ-14 & \vl{Validar el \textbf{Análisis y Diseño v1.1} (proceso to-be) como base del análisis con arquitectura.} & \vl{Documento} & \vl{Alta} & \vl{Presentado en reunión; sujeto al cierre de REQ-13} & \vl{En proceso} \\ \hline
\end{longtable}}

% ============================ 7. ACUERDOS =================================
\asec{7. ACUERDOS Y COMPROMISOS}
{\footnotesize\itshape\color{grisN} Cada compromiso debe expresar una acción verificable, un responsable único y una fecha acordada.\par}
\vspace{5pt}
{\footnotesize
\begin{longtable}{|C{0.8cm}|L{6.3cm}|L{2.5cm}|C{1.8cm}|C{1.5cm}|L{1.3cm}|}
\hline
\rowcolor{grisH}\hc{N°} & \hc{Acción / entregable} & \hc{Responsable} & \hc{Fecha límite} & \hc{Estado} & \hc{Evidencia} \\ \hline
\endfirsthead
\hline
\rowcolor{grisH}\hc{N°} & \hc{Acción / entregable} & \hc{Responsable} & \hc{Fecha límite} & \hc{Estado} & \hc{Evidencia} \\ \hline
\endhead
1 & \vl{Elaborar los \textbf{diagramas de casos de uso} y el \textbf{diagrama de BD con diccionario de datos} para cerrar el análisis} & \vl{Carlos / Luis (Analistas)} & \vl{Por definir (1--2 semanas tras confirmación)} & \vl{Pendiente} & \vl{Esta acta} \\ \hline
2 & \vl{\textbf{Decidir el modelo de usuarios} del portal: opción (a) 2 roles / 3 personas o (b) 3 roles internos} & \vl{Oscar (Cliente)} & \vl{08/09/2026 (en reunión)} & \vl{Cumplido} & \vl{Esta acta} \\ \hline
3 & \vl{Ajustar el RBAC y la gestión de usuarios según el modelo decidido} & \vl{Equipo UTEC (Backend)} & \vl{Fin de Sprint 1} & \vl{Pendiente} & \vl{Esta acta} \\ \hline
4 & \vl{Preparar el entorno de desarrollo (servidor + LLM de la universidad) para el Sprint 1} & \vl{Juan Ferreyra / Juan Renato} & \vl{11/09/2026} & \vl{En curso} & \vl{Plan de trabajo} \\ \hline
5 & \vl{Dejar \textbf{documentado} el alcance propuesto de 2ª fase (Cliente + tenant\_id, pasarela de pagos) como referencia futura --- \textbf{sin compromiso de desarrollo}: su ejecución dependerá del cierre de la fase 1 y de la decisión de presupuesto del cliente para salir a producción} & \vl{Fabricio (PM)} & \vl{09/09/2026} & \vl{En curso} & \vl{Correo} \\ \hline
6 & \vl{Enviar el acta 4 por correo dentro de las 24 horas} & \vl{Fabricio (PM)} & \vl{08/09/2026} & \vl{Pendiente} & \vl{Correo} \\ \hline
\end{longtable}}

% ============================ 8. PRÓXIMA ==================================
\asec{8. PRÓXIMA INTERACCIÓN Y SEGUIMIENTO}
{\footnotesize
\begin{tabular}{|L{2.5cm}|L{4.6cm}|L{2.6cm}|L{4.5cm}|}
\hline
\lb{Próxima reunión} & \vl{Por definir (según avance del análisis; plan refiere 11/09)} & \lb{Objetivo} & \vl{Confirmar diagrama de procesos/modelo de usuarios y revisar avance del análisis con arquitectura} \\ \hline
\lb{Canal de seguimiento} & \vl{Correo / Discord} & \lb{Responsable del acta} & \vl{Fabricio Godofredo Ladera La Torre (PM)} \\ \hline
\lb{Fecha de envío del acta} & \vl{08/09/2026} & \lb{Ubicación del archivo} & \vl{Repositorio del equipo CS3081} \\ \hline
\end{tabular}}

% ============================ 9. VALIDACIÓN ===============================
\asec{9. VALIDACIÓN DEL ACTA/FIRMAS}
{\footnotesize\itshape\color{grisN} La validación puede realizarse mediante firma, correo, mensaje en el canal acordado o confirmación registrada. Adjunte la evidencia correspondiente.\par}
\vspace{5pt}
{\footnotesize
\begin{tabular}{|L{4.7cm}|L{4.7cm}|L{4.7cm}|}
\hline
\hc{Representante del usuario / cliente} & \hc{Representante del equipo} & \hc{Docente o ACL (si corresponde)} \\ \hline
\vl{Nombre: Oscar Rafael Baldeón Montoro\newline Cargo / rol: Cliente / Product Owner} & \vl{Nombre: Fabricio Godofredo Ladera La Torre\newline Cargo / rol: Jefe de Proyecto (PM)} & \vl{Nombre: Teofilo Chambilla Aquino\newline Cargo / rol: Docente (acompañamiento)} \\ \hline
\vl{Método de validación: [Firma / correo]} & \vl{Método de validación: Firma} & \vl{Método de validación: [Firma / correo]} \\ \hline
\vl{Fecha: 08/09/2026\newline Observaciones: [Completar]} & \vl{Fecha: 08/09/2026\newline Observaciones: [Completar]} & \vl{Fecha: [dd/mm/aaaa]\newline Observaciones: [Completar]} \\ \hline
\end{tabular}}

% ============================ 10. APROBACIÓN ==============================
\asec{10. APROBACIÓN (V°B° CEO)}
El presente acta es revisada y aprobada por la Gerencia General antes de su envío al cliente.
\vspace{4pt}

{\itshape Oscar Rafael Baldeón Montoro}
\vspace{10pt}

\textbf{Nombre:} Oscar Rafael Baldeón Montoro

\textbf{Cargo:} CEO / Gerente General

\textbf{Firma:} \rule[-2pt]{6.5cm}{0.4pt}

\textbf{Fecha:} [dd/mm/aaaa]

\vspace{6pt}
\label{LastPage}
\end{document}
